\subsection{The ALTEN Group}
\label{sec:alten}

ALTEN is a French engineering and technology consulting group founded in 1988, with its
head office in Boulogne-Billancourt. It sells engineering expertise rather than products:
its consultants are placed on the technical projects of client companies, across the
aerospace, defence, energy, transport, telecommunications, health and finance sectors. The
group operates in more than thirty countries, and France remains its largest market.

Two features of that model matter for understanding this mission. First, ALTEN's asset is the
technical capability of its people, so anything that raises the level of a whole practice has
leverage that a single client project does not. Second, a consulting group works for many
clients at once, so it sees the same difficulty appear in company after company. That is how
a problem gets recognised as belonging to the industry rather than to one customer, and the
urgency assessment problem treated in this report is one of those.

Figure~\ref{fig:situation} places the mission inside the group, from the whole down to the
one layer this report is about. The rest of this section walks the same path.

\begin{figure}[H]
\centering
\begin{tikzpicture}[
    n/.style  = {bloc, minimum height=0.6cm, font=\scriptsize, align=left,
                 inner xsep=6pt},
    nF/.style = {blocFort, minimum height=0.6cm, font=\scriptsize, align=left,
                 inner xsep=6pt}]

\node[n, text width=11.2cm]  (g1)
  {\textbf{ALTEN Group:} engineering and technology consulting, 1988, thirty countries};
\node[n, text width=10.4cm, below right=1.8mm and 8mm of g1.south west,
      anchor=north west] (g2)
  {\textbf{Direction de l'Innovation:} the group's own R\&D, distinct from client work};
\node[n, text width=9.6cm, below right=1.8mm and 8mm of g2.south west,
      anchor=north west] (g3)
  {\textbf{Toulouse Lab:} one of the ALTEN Labs, where this mission took place};
\node[n, text width=8.8cm, below right=1.8mm and 8mm of g3.south west,
      anchor=north west] (g4)
  {\textbf{Smart Green Supply Chain:} lean, green and predictive aims at once};
\node[n, text width=8.0cm, below right=1.8mm and 8mm of g4.south west,
      anchor=north west] (g5)
  {\textbf{DISCO and its modules:} routing, supplier risk, carbon, transport};
\node[nF, text width=7.2cm, below right=1.8mm and 8mm of g5.south west,
      anchor=north west] (g6)
  {\textbf{This mission:} the urgency layer none of them provides};

\foreach \a/\b in {g1/g2, g2/g3, g3/g4, g4/g5, g5/g6}
  \draw[flecheFine] ($(\a.south west)+(4mm,0)$) |- (\b.west);

\end{tikzpicture}
\caption[Where the mission sits inside ALTEN]{From the group down to the layer this report is
about.}
\label{fig:situation}
\end{figure}

\subsection{The Direction de l'Innovation}
\label{sec:din}

The Direction de l'Innovation (DIN) is the entity responsible for ALTEN's own research and
development, as distinct from the work its consultants deliver to clients. It operates a
network of ALTEN Labs and brings together permanent researchers, doctoral candidates,
consultants, interns and apprentices. Its output is not billable hours but reusable assets:
methods, prototypes and expertise that consultants and clients can subsequently draw on.

This mission took place at the Toulouse Lab, within the Smart Green Supply Chain research
programme, which pursues three aims at once. The lean aim targets waste, meaning cost, delay
and excess stock. The green aim targets the carbon footprint of each transport leg, working to
the Global Logistics Emissions Council (GLEC) Framework
v3.2~\autocite{smartfreight2023glec} and ISO~14083:2023~\autocite{iso14083}. The predictive
aim, where this work sits, targets the anticipation of disruption.

The programme's central artefact is DISCO, a serious game and digital twin simulating a
complete supply chain for a satellite product across 22 component
suppliers~\autocite{baxas2025architecture}. Figure~\ref{fig:disco_map} shows the network
players work with and the three dimensions their decisions are scored on.

\begin{figure}[H]
    \centering
    \includegraphics[width=0.52\linewidth]{Images/1.Introduction/disco_map.png}
    \caption[The DISCO serious game network]{The DISCO serious game network. Players hold one
    node each and are scored on cost, delay and carbon.}
    \label{fig:disco_map}
\end{figure}

Successive internships have built its components: the Physical Internet routing
module~\autocite{baulain2024internet}, the supplier risk scorer~\autocite{kabouche2025scoring},
the per-route carbon quantification~\autocite{cordova2024green} and the transport management
interface~\autocite{mondelice2025tms}. The work described here adds a layer none of them
covers, and it ran alongside a parallel internship on the network representation and
geographic modelling side of the same programme, discussed in
Section~\ref{sec:parties_prenantes}. DISCO is the environment the present work was designed to
be compatible with, although the delivered system runs alongside it rather than inside it, for
reasons given in Section~\ref{sec:mission_perimetre}.

\subsection{Politique RSE d'ALTEN \textnormal{\textmd{(CSR policy, in French)}}}
\label{sec:rse}

\begin{frenchBox}
ALTEN a adh\'er\'e au Pacte Mondial des Nations Unies en 2010 et structure depuis sa politique
de responsabilit\'e soci\'etale autour de trois piliers~: l'humain, l'environnement et
l'innovation durable. La d\'emarche se mesure autant qu'elle se d\'eclare~: score EcoVadis de
85/100 en 2025, niveau Platinum et premier centile mondial, note A\textminus{} confirm\'ee au
Carbon Disclosure Project, et trajectoire de r\'eduction valid\'ee par la Science Based Targets
initiative en 2023 avec une neutralit\'e carbone vis\'ee \`a l'horizon 2050.

Sur le pilier humain, le Groupe r\'eunit plus de cent nationalit\'es, a sign\'e les
\emph{Women's Empowerment Principles} des Nations Unies et publie au titre de 2025 un index
d'\'egalit\'e professionnelle de 88/100 pour ALTEN Sud-Ouest, l'entit\'e juridique dans
laquelle s'est d\'eroul\'ee cette mission~; le programme INSPIRE 2026-2030 en porte la suite.
Sur le pilier environnemental, les actions internes sont concr\`etes et v\'erifiables~:
100~\% d'\'electricit\'e renouvelable pour les b\^atiments fran\c{c}ais dont la puissance
souscrite d\'epasse 36~kVA, 80~\% des surfaces couvertes par le tri s\'electif, une
certification ISO~14001 en cours d'extension, et des ateliers Fresque du Climat.

C'est le troisi\`eme pilier qui \'eclaire le mieux cette mission. ALTEN consacre 31~\% de sa
R\&D \`a l'innovation environnementale, et le programme dans lequel ce stage s'est inscrit en
fait partie~: un programme \emph{Smart Green} Supply Chain, o\`u les \'emissions de
CO\textsubscript{2} constituent l'un des six blocs du mod\`ele d'urgence d\'evelopp\'e ici, au
m\^eme rang que le co\^ut ou le d\'elai, et non une correction appliqu\'ee apr\`es coup. La
politique RSE du Groupe n'a donc pas seulement encadr\'e ce travail~: elle en a d\'etermin\'e
une variable de mod\'elisation.

\medskip

\noindent\textbf{Prise de recul.} Le carbone est l'un des six blocs du mod\`ele, et c'est
pourtant celui qui a le moins servi pendant la premi\`ere campagne. La raison n'est pas un
choix de mod\'elisation, c'est un manque de donn\'ees. Les \'emissions par trajet sont
connues \`a l'\'echelle du programme, pas encore \`a celle d'un n\oe{}ud et d'une semaine,
qui est la maille o\`u l'urgence se calcule. Une ambition environnementale peut donc \^etre
inscrite dans un mod\`ele et y rester la dimension la moins active. Avoir montr\'e \`a quel
endroit pr\'ecis la donn\'ee manque est, \`a mon sens, la contribution RSE la plus utile de
cette mission.
\end{frenchBox}
